\documentclass[11pt]{article}
\usepackage[a4paper,margin=0.9in]{geometry}
\usepackage[T1]{fontenc}
\usepackage{helvet}
\renewcommand{\familydefault}{\sfdefault}
\usepackage{xcolor}
\usepackage{booktabs}
\usepackage{array}
\usepackage{colortbl}
\usepackage{float}
\usepackage{amsmath}
\usepackage{titlesec}
\usepackage{fancyhdr}
\usepackage{enumitem}
\usepackage{lastpage}
\usepackage[hidelinks]{hyperref}

\definecolor{bcnavy}{RGB}{11,31,58}
\definecolor{bcblue}{RGB}{32,74,135}
\definecolor{bcgold}{RGB}{176,141,87}
\definecolor{bcgrey}{RGB}{90,98,112}
\definecolor{bcrule}{RGB}{210,214,220}
\definecolor{bckeep}{RGB}{232,237,244}
\definecolor{bcact}{RGB}{247,241,228}
\definecolor{bcwarn}{RGB}{250,238,238}

\titleformat{\section}{\color{bcnavy}\Large\bfseries}{\thesection}{0.6em}{}
\titleformat{\subsection}{\color{bcblue}\large\bfseries}{\thesubsection}{0.6em}{}
\titlespacing*{\section}{0pt}{1.3em}{0.45em}

\pagestyle{fancy}\fancyhf{}
\renewcommand{\headrulewidth}{0pt}
\renewcommand{\footrulewidth}{0.4pt}
\renewcommand{\footrule}{\hbox to\headwidth{\color{bcrule}\leaders\hrule height \footrulewidth\hfill}}
\fancyfoot[L]{\footnotesize\color{bcgrey}Bayesian Capital}
\fancyfoot[C]{\footnotesize\color{bcgrey}Confidential --- Internal Research}
\fancyfoot[R]{\footnotesize\color{bcgrey}Page \thepage\ of \pageref{LastPage}}

\setlength{\parindent}{0pt}\setlength{\parskip}{0.5em}

\newcommand{\note}[2]{%
  \begin{center}\fcolorbox{bcrule}{#1}{\parbox{0.95\textwidth}{\vspace{0.2em}#2\vspace{0.2em}}}\end{center}}

\begin{document}

{\color{bcgold}\rule{\textwidth}{1.6pt}}\\[0.3cm]
{\color{bcnavy}\huge\bfseries GM, VLO and CF}\\[0.15cm]
{\color{bcblue}\Large Diversifying the Daily Book Away from the AI Trade}\\[0.35cm]
{\color{bcgold}\rule{\textwidth}{0.8pt}}\\[0.35cm]

\begin{tabular}{@{}ll@{}}
\textbf{Date} & 6 August 2026 \\
\textbf{Candidates assessed} & GM, VLO, CF, ALNY \\
\textbf{Recommended} & GM, VLO, CF --- ALNY rejected \\
\textbf{Data} & 5-minute bars, 1 April 2024 to 3 August 2026, 587 sessions each \\
\textbf{Current book} & RKLB, TSM, VST, VRT, MU (daily model) \\
\textbf{Basis} & Verified fills, residual OU sigma, 50/50 sleeve split, marked to market \\
\end{tabular}

\vspace{0.6em}
{\color{bcrule}\rule{\textwidth}{0.6pt}}

\section{The question this answers}

Not ``do these names make money'' --- they make less than the book does, and that was never in
doubt. The question is whether they buy their way out of a concentration. Four of the five daily
names are one trade wearing different hats: TSM is semiconductors, VRT is data-centre cooling,
VST is power for data centres, MU is high-bandwidth memory. If the theme breaks they break
together, and a return cost is a rational price for not owning that.

The first pass scored these four on return and drawdown and rejected all of them. That answered
the wrong question and the conclusion is reversed here for three of them.

\section{The concentration, measured}

Measured on the \emph{shares}, not on the strategy. The strategy sits in cash much of the time,
which mutes every correlation and makes the book look better diversified than the assets are.

\begin{center}
\footnotesize
\begin{tabular}{@{}l r r r l@{}}
\toprule
& \textbf{Beta to AI factor} & \textbf{$R^2$} & \textbf{Beta to TSM} & \\
\midrule
RKLB & 0.77 & 0.18 & 0.69 & in the book \\
TSM  & 0.70 & 0.67 & 1.00 & in the book, AI \\
VST  & 1.00 & 0.63 & 0.77 & in the book, AI \\
VRT  & 1.19 & 0.78 & 1.06 & in the book, AI \\
MU   & 1.12 & 0.65 & 1.00 & in the book, AI \\
\midrule
GM   & 0.19 & 0.07 & 0.23 & candidate \\
VLO  & 0.11 & 0.02 & 0.06 & candidate \\
CF   & \textbf{0.00} & 0.00 & --0.04 & candidate \\
ALNY & 0.12 & 0.02 & 0.10 & candidate \\
\bottomrule
\end{tabular}
\end{center}

The AI factor is the equal-weight daily return of TSM, VRT, VST and MU --- the book's own
AI-adjacent names, no index being available. Beta to TSM alone is shown as a check that the
composite is not an artefact of how it was assembled; the two agree.

\note{bcwarn}{\textbf{RKLB is not the exception it looks like.} It is a space-launch company, not
an AI name, but its beta to the factor is 0.77 and to TSM 0.69. It does not share the theme, it
shares the risk appetite, and it sells off with it. The book is five-for-five exposed, not
four-for-five. This matters for sizing any replacement: dropping an AI name for RKLB would not
have helped.}

Average pairwise correlation of daily share returns: the four AI names \textbf{0.58}, the current
five \textbf{0.48}, the proposed nine \textbf{0.21}, the four candidates among themselves
\textbf{0.11}. One qualification --- VLO and CF correlate 0.41 with each other, both being
energy and commodity names. Treat them as one-and-a-bit names rather than two.

\section{Do the models work? The half-sample test}

Fit on the first half only, freeze, score the tested half. Boundary 23 May 2025. Nothing that
scores the tested half has seen it. Verified 5-minute fills throughout, which the earlier
diversifier shortlist could not have --- those names had no intraday coverage and were scored on
the at-open floor.

\begin{center}
\footnotesize
\begin{tabular}{@{}l r r r r r r@{}}
\toprule
& \textbf{Fit 1st half} & \textbf{TESTED half} & \textbf{Full fit} & \textbf{Buys/yr}
& \textbf{Tested DD} & \textbf{Corr to book} \\
\midrule
GM   & 68.3\% & \textbf{37.1\%} & 50.3\% & 94 & 10.8\% & 0.16 \\
VLO  & 29.6\% & \textbf{35.2\%} & 64.2\% & 20 &  4.3\% & 0.30 \\
CF   & 47.0\% & \textbf{42.4\%} & 46.9\% & 28 & 15.0\% & 0.01 \\
\rowcolor{bcwarn}
ALNY & 44.4\% & \textbf{--8.0\%} & 31.4\% & 23 & 35.8\% & 0.06 \\
\bottomrule
\end{tabular}
\end{center}

For contrast, the previous diversifier shortlist tested at --1.0\%, 4.1\%, 2.9\%, 15.1\% and
0.1\%. GM, VLO and CF are in a different class. VLO is the rarer case where the tested half beats
the fitted half, which is the opposite of the overfitting signature.

\textbf{ALNY fails here and is rejected.} 44.4\% fitted against --8.0\% tested, with a 35.8\%
drawdown, is the textbook pattern. Its share is as uncorrelated as the others, but a name whose
model does not work is not a diversifier --- it is a loss with low beta.

\section{A correction: the marginal test and rebalancing}

Whether a name belongs is not whether it clears a return bar. At equal weight a sixth name takes a
sixth of the capital, so it must beat what it dilutes.

The first version of this test built each book by averaging the names' \emph{daily} returns and
compounding --- equal weight restored every day. On that basis GM and CF looked close to free,
halving the book's drawdown at no cost. That construction is wrong: the book does not rebalance,
each sleeve owns its capital and compounds it independently, and nothing sweeps a winner's profit
back into a laggard overnight.

\begin{center}
\footnotesize
\begin{tabular}{@{}l r r r r r r@{}}
\toprule
& \multicolumn{3}{c}{\textbf{Rebalanced daily} (wrong basis)}
& \multicolumn{3}{c}{\textbf{Held} (as the book runs)} \\
\cmidrule(lr){2-4}\cmidrule(lr){5-7}
\textbf{Tested half} & Return & Max DD & Sharpe & Return & Max DD & Sharpe \\
\midrule
five names & 111.1\% & 29.7\% & 2.22 & \textbf{127.4\%} & \textbf{14.9\%} & \textbf{2.75} \\
+ GM       & 109.8\% & 16.3\% & 2.85 & 113.0\% & 13.3\% & 2.78 \\
+ VLO      & 108.6\% & 17.4\% & 2.88 & 112.6\% & 13.7\% & 2.80 \\
+ CF       & 110.5\% & 14.9\% & 2.94 & 113.7\% & 13.0\% & 2.83 \\
+ ALNY     &  90.5\% & 22.2\% & 2.58 & 106.3\% & 14.4\% & 2.68 \\
\bottomrule
\end{tabular}
\end{center}

Corrected, the book's drawdown is already 14.9\%, not 29.7\% --- the 29.7\% \emph{was} the
rebalancing, feeding capital back into names while they fell. Rebalancing is a drag here, not a
bonus, costing the five-name book 16.3pp: with strongly trending names it sells the winners.

On this basis each candidate costs 13 to 21pp of return to remove 1 to 2pp of drawdown, and
Sharpe barely moves. \textbf{On a return-and-drawdown objective, all four should be rejected.}
That is the correct conclusion to that question, and it is not the question being asked.

\section{The stress windows: what the concentration actually costs}

Concentration is a claim about bad days, so it is tested on bad days. The deepest drawdown
windows of the AI factor were located automatically rather than chosen.

\begin{center}
\footnotesize
\begin{tabular}{@{}l r r r@{}}
\toprule
& \textbf{Jan--Apr 2025} & \textbf{Jun--Jul 2026} & \textbf{May--Aug 2024} \\
& \tiny AI factor --45.9\% & \tiny AI factor --17.6\% & \tiny AI factor --27.3\% \\
\midrule
five names           & \textbf{--37.3\%} & --9.2\% & --4.5\% \\
nine names           & \textbf{--18.9\%} & --2.7\% & --1.8\% \\
five + GM, VLO, CF   & --21.6\%          & --2.9\% & --2.8\% \\
\bottomrule
\end{tabular}
\end{center}

In the January to April 2025 episode every book share fell together --- VRT --60.0\%,
VST --47.5\%, RKLB --44.7\%, MU --40.7\%, TSM --34.2\% --- while the candidates lost far less:
GM --16.2\%, CF --21.8\%, VLO --22.5\%, ALNY --11.7\%. In June to July 2026, VLO rose 30.9\% and
CF rose 17.9\% \emph{while} the theme fell. That is the hedge; the OU sleeve never delivered one.

Over the full sample, adding the three takes the book's beta to the factor from 0.58 to 0.41 and
its maximum drawdown from 39.3\% to 27.0\%.

\note{bckeep}{\textbf{One nuance that qualifies the win.} $R^2$ to the AI factor stays at 0.37 for
the five-name book, the eight and the nine alike, while beta falls from 0.58 to 0.41. Adding these
names scales total book risk down by about a third rather than changing what the risk \emph{is}.
Beta and the stress windows are the honest measures of reliance here; $R^2$ is not.}

\section{Walk-forward, and what it cannot show}

Three expanding folds, every name --- the five incumbents included --- re-fitted on each training
window, no rebalancing. Re-fitting the incumbents matters: their deployed parameters were fitted
on the whole sample, so scoring those against candidates frozen from train would let the incumbent
see its own test data.

\begin{center}
\footnotesize
\begin{tabular}{@{}l l r r r r r r@{}}
\toprule
\textbf{Fold} & \textbf{Test window} & \textbf{AI factor} & \textbf{Five} & \textbf{Eight}
& \textbf{Delta} & \textbf{DD five} & \textbf{DD eight} \\
\midrule
1 & 2025-05-28 to 2025-10-13 & +61.2\% & 120.0\% & 79.1\% & --41.0 &  4.6\% & 2.8\% \\
2 & 2025-10-14 to 2026-03-04 & +31.1\% &  69.2\% & 56.7\% & --12.5 & 11.4\% & 8.3\% \\
3 & 2026-03-05 to 2026-07-23 & +42.5\% &  49.8\% & 61.6\% & \textbf{+11.8} & 16.1\% & 8.9\% \\
\bottomrule
\end{tabular}
\end{center}

The eight-name book has a lower drawdown in 3 of 3 folds, mean --4.0pp, and a lower return in
2 of 3, mean --13.9pp. The cost tracks how hard the theme ran, as a diversifier's should: 41pp
given up when the factor rose 61\%, 12.5pp when it rose 31\%.

\textbf{Fold 3 is the most informative.} The eight-name book won by 11.8pp because RKLB fell
28.6\% while VLO rose 50.6\% and CF rose 20.2\%. A book name broke for its own reasons and the
diversifiers carried it --- the thesis in miniature, with the theme itself still up.

\note{bcwarn}{\textbf{What the walk-forward does not test.} The AI factor \emph{rose} in every
test window: +61.2\%, +31.1\%, +42.5\%. The only AI drawdown in the sample, January to April 2025,
falls before the first fold boundary and therefore sits in the training half of all three folds.
The walk-forward prices the premium and never tests the insurance. The halved stress loss in
Section 5 remains evidence from a single episode inside the fitted region. A reversed
walk-forward --- fit late, score the crash --- would put that episode out of sample, at the cost
of inverting time order.}

\section{Planning returns}

Taken only from windows nothing had seen: the three walk-forward folds, plus for these names the
half-sample fit-then-freeze test, which is the observation closest to how they would be run.
Median rather than mean, so a single fold does not set the number.

\begin{center}
\footnotesize
\begin{tabular}{@{}l r r r r r@{}}
\toprule
& \textbf{Fold 1} & \textbf{Fold 2} & \textbf{Fold 3} & \textbf{Half-sample}
& \textbf{Planning} \\
\midrule
GM  & 38.2\% & 28.9\% &  29.5\% & 37.1\% & \textbf{33\%} \\
VLO & 35.0\% & 34.4\% & 191.0\% & 35.2\% & \textbf{35\%} \\
CF  &  1.6\% & 40.4\% &  61.6\% & 42.4\% & \textbf{41\%} \\
\bottomrule
\end{tabular}
\end{center}

The two routes were derived independently and agree: the fold medians land within 4pp of the
half-sample figures. GM's four observations span 9pp and VLO has three within a point of each
other, steadier than any book name manages.

\textbf{A basis check, since the book's figures are quoted differently.} The book's published
planning figures come from full-sample fits; these come from frozen ones. Scoring the five book
names on the same unseen-window basis gives a mean of 77.0\% against the published 79.4\% --- about
2pp. The two can be quoted side by side without correction.

\note{bckeep}{\textbf{Plan on 35\% for each of the three}, with CF slightly above and GM slightly
below; the difference is well inside the noise. For the Allocation sheet's return-assumption
column, 0.35 for all three. \\[0.4em]
\textbf{Cost to the book: about 79\% to about 64\%, some 15pp.} That is an independent route to
the walk-forward's --13.9pp, which is reassuring. On an equal-weight eight-name book each name
takes 12.5\% of capital.}

\textbf{Do not read the tightness as precision.} Annualising a four-and-a-half month window
amplifies everything. The book names make the point: on the same arithmetic MU's fold median comes
out at 211\% and RKLB's at --1.8\%, both useless as planning figures. The new names' numbers are
tighter because their models trade steadier, but it is the same calculation. CF's fold 1 was
1.6\% --- the 41\% is a median, not a floor.

\section{Parameters for the model}

Fitted on the whole sample with the robust differential-evolution objective, verified fills,
residual OU sigma, 50/50 split --- the same procedure and the same basis the book's own names
were fitted on. Full-sample figures: GM 50.3\%, VLO 64.2\%, CF 46.9\%, at 61, 35 and 27 buys a
year respectively. Those are in-sample and are not the planning figures; Section 7 is.

\vspace{0.4em}
\begin{center}
\footnotesize
\begin{tabular}{@{}l r r r r r r@{}}
\toprule
& \multicolumn{2}{c}{\textbf{GM}} & \multicolumn{2}{c}{\textbf{VLO}}
& \multicolumn{2}{c}{\textbf{CF}} \\
\cmidrule(lr){2-3}\cmidrule(lr){4-5}\cmidrule(lr){6-7}
& \textbf{deploy} & \tiny 1st half & \textbf{deploy} & \tiny 1st half
& \textbf{deploy} & \tiny 1st half \\
\midrule
$\lambda$      & \textbf{0.704472}   & 0.2011   & \textbf{0.723393}  & 1.441    & \textbf{0.736218}  & 0.8737 \\
$\varphi_L$    & \textbf{0.750723}   & 0.2907   & \textbf{0.625222}  & 0.7016   & \textbf{0.106120}  & 0.4777 \\
$\psi$         & \textbf{0.0805582}  & 0.04228  & \textbf{0.0200999} & 0.02011  & \textbf{0.0694359} & 0.004598 \\
$k$            & \textbf{0.461711}   & 0.3827   & \textbf{0.543763}  & 0.9450   & \textbf{2.03552}   & 1.838 \\
Premium        & \textbf{0.00969721} & 0.005803 & \textbf{0.0303035} & 0.02321  & \textbf{0.0465592} & 0.02545 \\
Peak cap       & \textbf{0.0423717}  & 0.03368  & \textbf{0.0308569} & 0.05922  & \textbf{0.0286416} & 0.06327 \\
OU buffer      & \textbf{1.06024}    & 0.6041   & \textbf{1.17841}   & 1.799    & \textbf{0.841237}  & 0.8036 \\
OU premium     & \textbf{0.0451112}  & 0.009885 & \textbf{0.0419514} & 0.02552  & \textbf{0.0590932} & 0.05159 \\
OU cap         & \textbf{0.0417662}  & 0.09131  & \textbf{0.0370821} & 0.08021  & \textbf{0.0831024} & 0.09228 \\
OU lookback $W$& \textbf{45}         & 41       & \textbf{125}       & 117      & \textbf{66}        & 66 \\
\bottomrule
\end{tabular}
\end{center}

\note{bcwarn}{\textbf{The fits move a great deal between halves, and the returns do not.} GM's
$\lambda$ goes from 0.20 to 0.70 and its OU premium from 0.0099 to 0.045; CF's $\psi$ moves
fifteen-fold. Yet the out-of-sample returns those two parameter sets produce sit within a few
points of each other. The optimum is broad rather than sharp, which is reassuring about the
result and a warning about the vectors: do not treat these numbers as precise, and do not re-fit
them expecting an improvement. The one parameter that holds still is the OU lookback --- 41 to 45,
117 to 125, 65.6 to 65.6 --- which is the structural one.}

Shared with the rest of the book: commission \$0.005 per share, IBKR interest 3.14\% p.a.,
stop-loss 50 calendar days, Bayes share 0.50, residual OU sigma. Note the OU buffer is on the
residual scale and is not interchangeable with a level-sigma buffer.

\section{Recommendation}

\note{bcact}{\textbf{Add GM, VLO and CF. Reject ALNY.} Plan each of the three at 35\%, taking the
book from roughly 79\% to roughly 64\%. In exchange, beta to the AI factor falls from 0.58 to 0.41,
full-sample maximum drawdown from 39.3\% to 27.0\%, and the loss through the one AI drawdown in
the sample from 37.3\% to 21.6\%. ALNY is rejected on its model, not its correlation: --8.0\% on
the tested half against 44.4\% fitted.}

What should temper it:

\begin{itemize}[leftmargin=1.4em,itemsep=0.15em]
\item The stress evidence rests on \textbf{one episode}, and that episode is inside the region
      the parameters were fitted on. No unseen window in this study contains a falling AI market.
\item \textbf{VLO and CF correlate 0.41.} Funding both buys less independent diversification than
      the name count suggests.
\item Everything here comes from \textbf{2.3 years of a broadly rising market}. VLO and CF are
      energy and commodity names whose behaviour over this window may not represent their regime.
\item \textbf{The three add 123 buys a year} --- GM 61, VLO 35, CF 27 --- against the current
      book's 335 across five names, so about 37\% more order flow for names that dilute return.
\item These three are not credible as \textbf{replacements}. TSM is the weakest book name at
      82.4\% on the tested half and still doubles the best of them.
\end{itemize}

\vspace{0.8em}
{\color{bcrule}\rule{\textwidth}{0.6pt}}

{\footnotesize\color{bcgrey}
All figures verified against 5-minute bars, marked to market, net of \$0.005 per share commission
and inclusive of interest on idle cash at 3.14\%. Books are constructed without rebalancing
between names, matching how the sleeves run. Past backtested performance is not a forecast of
future returns.}

\end{document}
